% !TeX program = xelatex
% UTF-8; compile twice with XeLaTeX. No external images or data files needed.
\documentclass[a4paper,landscape,10pt]{article}
\usepackage[margin=14mm,headheight=14pt,headsep=5mm,footskip=7mm]{geometry}
\usepackage{fontspec}
\setmainfont{DejaVu Sans}
\newfontfamily\ipafont{DejaVu Sans}
\usepackage[table]{xcolor}
\usepackage{array,longtable,booktabs,fancyhdr}
\definecolor{ink}{HTML}{18384D}
\definecolor{band}{HTML}{E7EFF4}
\pagestyle{fancy}\fancyhf{}
\fancyhead[L]{\small\color{ink}Parlez Vous QSO?! · Français / French}
\fancyhead[R]{\small Englische Phrase · Traduction française · IPA}
\fancyfoot[L]{\scriptsize Vorlage: French v 1.0 · Credits: Jean-Michael -- F5OQO; Pierre-Louis -- F5BTT}
\fancyfoot[R]{\small\thepage}
\setlength{\parindent}{0pt}
\setlength{\tabcolsep}{3mm}
\renewcommand{\arraystretch}{1.25}
\newcommand{\ipa}[1]{{\ipafont[#1]}}
\begin{document}
{\LARGE\bfseries\color{ink}Parlez Vous QSO?!}\hfill{\large Français / French}

\vspace{2mm}
{\small QSO-Sprachführer mit französischer IPA-Lautschrift. Die Lautschrift bezieht sich auf die französische Spalte.}

\vspace{2mm}
{\footnotesize \textbf{Leseschlüssel:} \ipa{ɑ̃, ɛ̃, ɔ̃, œ̃} Nasalvokale; \ipa{y} wie deutsches ü,
\ipa{ɥ} kurzer ü-Gleitlaut, \ipa{ʁ} französisches r. \ipa{‿} verbindet Wörter in der Aussprache.
Schrägstriche trennen Alternativen; \ipa{…} steht für einen einzusetzenden Ausdruck.}

\vspace{1mm}
\fontsize{9}{11.8}\selectfont
\begin{longtable}{@{}>{\raggedright\arraybackslash}p{13mm}>{\raggedright\arraybackslash}p{73mm}>{\raggedright\arraybackslash}p{75mm}>{\raggedright\arraybackslash}p{90mm}@{}}
\toprule
\textbf{Ref.} & \textbf{English phrase} & \textbf{Français} & \textbf{Lautschrift (IPA)}\\
\midrule
\endfirsthead
\toprule
\textbf{Ref.} & \textbf{English phrase} & \textbf{Français} & \textbf{Lautschrift (IPA)}\\
\midrule
\endhead
\midrule
\multicolumn{4}{r}{\scriptsize Fortsetzung auf der nächsten Seite}\\
\endfoot
\bottomrule
\endlastfoot
\rowcolor{band} 1.0 & \textbf{Opening Greetings} & \multicolumn{2}{l}{\textbf{Début de QSO}} \\*
1.1 & I am very pleased to meet you. & Je suis très heureux de vous contacter. & \ipa{ʒə sɥi tʁɛ.z‿øʁø də vu kɔ̃takte} \\
1.2 & Thank you for your call & Merci de votre appel. & \ipa{mɛʁsi də vɔtʁ apɛl} \\
1.3 & Good morning & Bonjour. & \ipa{bɔ̃ʒuʁ} \\
1.31 & Good day & Bonjour. & \ipa{bɔ̃ʒuʁ} \\
1.32 & Good afternoon & Bonjour. & \ipa{bɔ̃ʒuʁ} \\
1.33 & Good evening & Bonsoir. & \ipa{bɔ̃swaʁ} \\
1.4 & I know only a little French but I enjoy QSOs in French. & Je connais peu de mots en français mais j'aime les QSO en français. & \ipa{ʒə kɔnɛ pø də mo.z‿ɑ̃ fʁɑ̃sɛ mɛ ʒɛm le ky ɛs o ɑ̃ fʁɑ̃sɛ} \\
\rowcolor{band} 2.0 & \textbf{Operator} & \multicolumn{2}{l}{\textbf{Opérateur}} \\*
2.1 & My name is Peter. & Mon nom est Pierre. & \ipa{mɔ̃ nɔ̃ ɛ pjɛʁ} \\
2.2 & I have been licensed since 1960. & J’ai ma licence depuis 1930. & \ipa{ʒe ma lisɑ̃s dəpɥi mil nœf sɑ̃ tʁɑ̃t} \\
2.3 & I have been licensed for 43 years. & J’ai mon indicatif depuis 43 ans. & \ipa{ʒe mɔ̃.n‿ɛ̃dikatif dəpɥi kaʁɑ̃t tʁwa.z‿ɑ̃} \\
2.4 & I live in Sittingbourne. & J'habite Sittingbourne. & \ipa{ʒabit ˈsɪtɪŋbɔːn} \\
2.5 & Sittingbourne is a village. & Sittingbourne est un village. & \ipa{ˈsɪtɪŋbɔːn ɛ.t‿œ̃ vilaʒ} \\
2.51 & Sittingbourne is a town. & Sittingbourne est une petite ville. & \ipa{ˈsɪtɪŋbɔːn ɛ.t‿yn pətit vil} \\
2.52 & Sittingbourne is a city. & Sittingbourne est une grande ville. & \ipa{ˈsɪtɪŋbɔːn ɛ.t‿yn ɡʁɑ̃d vil} \\
2.6 & Sittingbourne is between London and Dover. & Sittingbourne est située entre Londres et Douvres. & \ipa{ˈsɪtɪŋbɔːn ɛ sitɥe ɑ̃tʁ lɔ̃dʁ e duvʁ} \\
2.7 & I live in the South of England. & J'habite au sud de l’Angleterre. & \ipa{ʒabit o syd də lɑ̃ɡlətɛʁ} \\
2.71 & I live in the North of England. & J'habite au nord de l’Angleterre. & \ipa{ʒabit o nɔʁ də lɑ̃ɡlətɛʁ} \\
2.72 & I live in the East of England. & J'habite à l’est de l’Angleterre. & \ipa{ʒabit a lɛst də lɑ̃ɡlətɛʁ} \\
2.73 & I live in West of England. & J'habite à l’ouest de l’Angleterre. & \ipa{ʒabit a lwɛst də lɑ̃ɡlətɛʁ} \\
2.74 & I live in South East England. & J'habite au sud est de l’Angleterre. & \ipa{ʒabit o syd ɛst də lɑ̃ɡlətɛʁ} \\
2.75 & I live in South West England. & J'habite au sud ouest de l’Angleterre. & \ipa{ʒabit o syd wɛst də lɑ̃ɡlətɛʁ} \\
2.76 & I live in North East England. & J'habite au nord est de l’Angleterre & \ipa{ʒabit o nɔʁ ɛst də lɑ̃ɡlətɛʁ} \\
2.77 & I live in North West England. & J'habite au nord ouest de l’Angleterre. & \ipa{ʒabit o nɔʁ wɛst də lɑ̃ɡlətɛʁ} \\
2.78 & I live in central England. & J'habite au centre de l’Angleterre. & \ipa{ʒabit o sɑ̃tʁ də lɑ̃ɡlətɛʁ} \\
2.8 & I live in the country. & J'habite à la campagne. & \ipa{ʒabit a la kɑ̃paɲ} \\
2.81 & I live by the sea. & J'habite près de la mer. & \ipa{ʒabit pʁɛ də la mɛʁ} \\
2.82 & I live near mountains. & J'habite à la montagne. & \ipa{ʒabit a la mɔ̃taɲ} \\
2.9 & Sittingbourne is 50Km from London. & Sittingbourne est à 50 Km de Londres. & \ipa{ˈsɪtɪŋbɔːn ɛ.t‿a sɛ̃kɑ̃t kilɔmɛtʁ də lɔ̃dʁ} \\
2.11 & I am retired. & Je suis à la retraite.\newline Je suis retraité. & \ipa{ʒə sɥi.z‿a la ʁətʁɛt / ʒə sɥi ʁətʁɛte} \\
2.12 & My home call sign is M0ECS. & Mon indicatif est M0ECS. & \ipa{mɔ̃.n‿ɛ̃dikatif ɛ ɛm zeʁo ø se ɛs} \\
\rowcolor{band} 3.0 & \textbf{Station} & \multicolumn{2}{l}{\textbf{Station}} \\*
3.1 & My radio is an Icom 706. & Mon transceiver est un Icom 706. & \ipa{mɔ̃ tʁɑ̃sivœʁ ɛ.t‿œ̃ ikɔm sɛt sɑ̃ sis} \\
3.2 & My station is all homebrew (homemade). & Ma station est entièrement de fabrication personnelle. & \ipa{ma stasjɔ̃ ɛ.t‿ɑ̃tjɛʁmɑ̃ də fabʁikasjɔ̃ pɛʁsɔnɛl} \\
3.3 & I am running 5 Watts. & Ma puissance est de 5 watts. & \ipa{ma pɥisɑ̃s ɛ də sɛ̃k wat} \\
3.4 & My antenna is a long wire. & Mon antenne est un long fil. & \ipa{mɔ̃.n‿ɑ̃tɛn ɛ.t‿œ̃ lɔ̃ fil} \\
3.41 & My antenna is a dipole. & Mon antenne est un dipôle. & \ipa{mɔ̃.n‿ɑ̃tɛn ɛ.t‿œ̃ dipol} \\
3.42 & My antenna is a G5RV. & Mon antenne est une G5RV. & \ipa{mɔ̃.n‿ɑ̃tɛn ɛ.t‿yn ʒe sɛ̃k ɛʁ ve} \\
3.43 & My antenna is a doublet. & Mon antenne est un doublet. & \ipa{mɔ̃.n‿ɑ̃tɛn ɛ.t‿œ̃ dublɛ} \\
3.44 & My antenna is a delta loop. & Mon antenne est une delta loop. & \ipa{mɔ̃.n‿ɑ̃tɛn ɛ.t‿yn dɛlta lup} \\
3.45 & My antenna is a vertical. & Mon antenne est une verticale. & \ipa{mɔ̃.n‿ɑ̃tɛn ɛ.t‿yn vɛʁtikal} \\
3.46 & My antenna is a log periodic. & Mon antenne est une log périodique. & \ipa{mɔ̃.n‿ɑ̃tɛn ɛ.t‿yn lɔɡ peʁjɔdik} \\
3.47 & My antenna is a Yagi beam. & Mon antenne est une Yagi & \ipa{mɔ̃.n‿ɑ̃tɛn ɛ.t‿yn jaɡi} \\
3.48 & My antenna is a 5 element beam. & Mon antenne est une Yagi (de) 5 éléments. & \ipa{mɔ̃.n‿ɑ̃tɛn ɛ.t‿yn jaɡi (də) sɛ̃k e.le.mɑ̃} \\
\rowcolor{band} 4.0 & \textbf{Conditions} & \multicolumn{2}{l}{\textbf{Conditions}} \\*
4.1 & Thank you for my 599 report. & Merci pour mon RST de 599. & \ipa{mɛʁsi puʁ mɔ̃.n‿ɛʁ ɛs te də sɛ̃k nœf nœf} \\
4.2 & The band conditions are good. & Les conditions sur la bande sont bonnes & \ipa{le kɔ̃disjɔ̃ syʁ la bɑ̃d sɔ̃ bɔn} \\
4.3 & The band conditions are bad. & Les conditions sur la bande sont mauvaises. & \ipa{le kɔ̃disjɔ̃ syʁ la bɑ̃d sɔ̃ movɛz} \\
4.4 & There is strong interference here. & Il y a de fortes interférences ici. & \ipa{il j a də fɔʁt.z‿ɛ̃tɛʁfeʁɑ̃s isi} \\
4.5 & Can you change to the 80m band? & Pouvez vous venir sur la bande des 80 mètres. & \ipa{puve vu vəniʁ syʁ la bɑ̃d de katʁə.vɛ̃ mɛtʁ} \\
4.6 & Please change frequency to 3.558 MHz. & Réglez vous sur 3,556 MHz s'il vous plaît. & \ipa{ʁeɡle vu syʁ tʁwa viʁɡyl sɛ̃k sɛ̃k sis meɡaɛʁts sil vu plɛ} \\
\rowcolor{band} 5.0 & \textbf{General Queries and Comments} & \multicolumn{2}{l}{\textbf{Questions Générales et Commentaires}} \\*
5.1 & What does <\#\#\#\#\#> mean please? & Que signifie <\#\#\#\#\#> s'il vous plaît ? & \ipa{kə siɲifi … sil vu plɛ} \\
5.2 & Please wait a moment, I need to check my French dictionary. & Un instant s'il vous plaît, je vérifie dans mon dictionnaire français. & \ipa{œ̃.n‿ɛ̃stɑ̃ sil vu plɛ ʒə veʁifi dɑ̃ mɔ̃ diksjɔnɛʁ fʁɑ̃sɛ} \\
5.3 & How do you say, <\#\#\#\#\#> in French ? & Comment dit-on <\#\#\#\#\#> en français? & \ipa{kɔmɑ̃ di t‿ɔ̃ … ɑ̃ fʁɑ̃sɛ} \\
5.5 & My favourite modes of operation are CW / RTTY / PSK / MT63 / Olivia / Hellschreiber / Digital Modes / and SSB & Mes modes préférés sont CW / RTTY / PSK / MT63 / Olivia / Hellschreiber / modes digitaux / et SSB. & \ipa{me mɔd pʁefeʁe sɔ̃ se dublə ve / ɛʁ te te i ɡʁɛk / pe ɛs ka / ɛm te swasɑ̃t tʁwa / ɔlivja / ˈhɛlˌʃʁaɪbɐ / mɔd diʒito / e ɛs ɛs be} \\
5.6 & My best DX is Australia. & Mon meilleur DX est l'Australie. & \ipa{mɔ̃ mɛjœʁ de iks ɛ lostʁali} \\
5.7 & I enjoy working contests. & J'aime les concours. & \ipa{ʒɛm le kɔ̃kuʁ} \\
5.8 & I like working QRP. & J’aime trafiquer en QRP. & \ipa{ʒɛm tʁafike ɑ̃ ky ɛʁ pe} \\
\rowcolor{band} 6.0 & \textbf{The Weather} & \multicolumn{2}{l}{\textbf{Le Temps}} \\*
6.1 & It is sunny. & Il y a du soleil. & \ipa{il j a dy sɔlɛj} \\
6.2 & It is windy. & Il y a du vent. & \ipa{il j a dy vɑ̃} \\
6.3 & It is raining. & Il pleut. & \ipa{il plø} \\
6.4 & It is snowing. & Il neige. & \ipa{il nɛʒ} \\
6.5 & It is very hot. & Il fait très chaud. & \ipa{il fɛ tʁɛ ʃo} \\
6.6 & It is very cold. & Il fait très froid. & \ipa{il fɛ tʁɛ fʁwa} \\
6.7 & The temperature is 8C. & La température est de 8°C. & \ipa{la tɑ̃peʁatyʁ ɛ də ɥit dəɡʁe sɛlsjys} \\
6.8 & It is overcast. & Le temps est couvert. & \ipa{lə tɑ̃.z‿ɛ kuvɛʁ} \\
6.9 & It is foggy. & Le temps est brumeux. & \ipa{lə tɑ̃.z‿ɛ bʁymø} \\
6.11 & It has been raining all day. & Il a plu toute la journée. & \ipa{il a ply tut la ʒuʁne} \\
6.12 & There is 20cm of snow here. & Il y a 20 cm de neige ici. & \ipa{il j a vɛ̃ sɑ̃timɛtʁ də nɛʒ isi} \\
\rowcolor{band} 8.0 & \textbf{QSL} & \multicolumn{2}{l}{\textbf{QSL}} \\*
8.1 & My QSL card is sure via the bureau & Ma carte QSL est sûre via bureau. & \ipa{ma kaʁt ky ɛs ɛl ɛ syʁ vja byʁo} \\
8.2 & QSL via eQSL & QSL via eQSL. & \ipa{ky ɛs ɛl vja i ky ɛs ɛl} \\
8.3 & No eQSL please & Pas de eQSL s'il vous plaît. & \ipa{pa də i ky ɛs ɛl sil vu plɛ} \\
8.4 & I will send you my card direct. & Je vous envoie ma carte QSL en direct. & \ipa{ʒə vu.z‿ɑ̃vwa ma kaʁt ky ɛs ɛl ɑ̃ diʁɛkt} \\
8.5 & What is your address? & Quelle est votre adresse ? & \ipa{kɛl ɛ vɔtʁ adʁɛs} \\
8.6 & Please send your QSL card direct. & Envoyez (moi) votre carte QSL en direct s'il vous plaît. & \ipa{ɑ̃vwaje (mwa) vɔtʁ kaʁt ky ɛs ɛl ɑ̃ diʁɛkt sil vu plɛ} \\
8.7 & My address is... & Mon adresse est ... & \ipa{mɔ̃.n‿adʁɛs ɛ …} \\
8.8 & My address is shown on QSL.com / EHam.net & Mon adresse est indiquée sur QSL.com / EHam.net. & \ipa{mɔ̃.n‿adʁɛs ɛ.t‿ɛ̃dike syʁ ky ɛs ɛl pwɛ̃ kɔm / i hæm pwɛ̃ nɛt} \\
8.9 & Sorry I do not send QSL cards. & Désolé je n'envoie pas de carte QSL. & \ipa{dezɔle ʒə nɑ̃vwa pa də kaʁt ky ɛs ɛl} \\
8.11 & Please send QSL card via G1JQH. & Envoyez les cartes QSL via G1JQH s'il vous plaît. & \ipa{ɑ̃vwaje le kaʁt ky ɛs ɛl vja ʒe œ̃ ʒi ky aʃ sil vu plɛ} \\
8.12 & My WACRAL / GQRP / DXCC / FISTS number is 1234. & Mon numéro WACRAL / GQRP / DXCC / FISTS est 1234. & \ipa{mɔ̃ nymeʁo dublə ve a se ɛʁ ɛl / ʒe ky ɛʁ pe / de iks se se / fɪsts ɛ mil dø sɑ̃ tʁɑ̃t katʁ} \\
\rowcolor{band} 9.0 & \textbf{Final Greetings} & \multicolumn{2}{l}{\textbf{Fin de QSO}} \\*
9.1 & Thank you. & Merci. & \ipa{mɛʁsi} \\
9.2 & Thank you for a very enjoyable contact. & Merci pour ce contact très sympathique. & \ipa{mɛʁsi puʁ sə kɔ̃takt tʁɛ sɛ̃patik} \\
9.3 & I enjoyed our chat. & J'ai apprécié notre discussion. & \ipa{ʒe apʁesje nɔtʁ diskysjɔ̃} \\
9.4 & I enjoyed our QSO in French. & J'ai apprécié notre QSO en français. & \ipa{ʒe apʁesje nɔtʁ ky ɛs o ɑ̃ fʁɑ̃sɛ} \\
9.5 & It has been great fun. & Cela a été très intéressant. & \ipa{səla a ete tʁɛ.z‿ɛ̃teʁesɑ̃} \\
9.6 & If you would like a regular contact with me, you can contact me via email. & Si vous voulez me contacter régulièrement, vous pouvez le faire par email. & \ipa{si vu vule mə kɔ̃takte ʁeɡyljɛʁmɑ̃ vu puve lə fɛʁ paʁ imɛl} \\
9.7 & My email address is... & Mon adresse email est ... & \ipa{mɔ̃.n‿adʁɛs imɛl ɛ …} \\
9.8 & My email address is shown on QSL.com / EHam.net & Mon adresse email est sur QSL.com / EHam.net. & \ipa{mɔ̃.n‿adʁɛs imɛl ɛ syʁ ky ɛs ɛl pwɛ̃ kɔm / i hæm pwɛ̃ nɛt} \\
9.9 & I look forward to meeting you again. & J'espère vous retrouver bientôt & \ipa{ʒɛspɛʁ vu ʁətʁuve bjɛ̃to} \\
9.11 & My best wishes to you and your family. & Meilleurs souhaits ( / vœux) pour vous et votre famille. & \ipa{mɛjœʁ suɛ / mɛjœʁ vø puʁ vu e vɔtʁ famij} \\
9.12 & Best wishes to your wife. & Meilleurs souvenirs à votre épouse. & \ipa{mɛjœʁ suvəniʁ a vɔtʁ epuz} \\
9.121 & Best wishes to your husband. & Meilleurs souhaits à votre époux. & \ipa{mɛjœʁ suɛ.z‿a vɔtʁ epu} \\
9.13 & Good luck. & Bonne chance. & \ipa{bɔn ʃɑ̃s} \\
9.14 & Good health. & Bonne santé. & \ipa{bɔn sɑ̃te} \\
9.15 & Goodbye & Au revoir. & \ipa{o ʁəvwaʁ} \\
9.16 & I wish you an enjoyable afternoon. & Je vous souhaite un bon après midi. & \ipa{ʒə vu suɛt œ̃ bɔ̃.n‿apʁɛ midi} \\
9.161 & I wish you an enjoyable evening. & Je vous souhaite une bonne soirée. & \ipa{ʒə vu suɛt yn bɔn swaʁe} \\
9.162 & I wish you an enjoyable weekend. & Je vous souhaite un bon week-end. & \ipa{ʒə vu suɛt œ̃ bɔ̃ wikɛnd} \\
9.17 & Enjoy your holiday. & Bonnes vacances. & \ipa{bɔn vakɑ̃s} \\
9.18 & Happy New Year. & Bonne année.\newline Bonne année, bonne santé & \ipa{bɔ.n‿a.ne / bɔ.n‿a.ne bɔn sɑ̃te} \\
\rowcolor{band} 10.0 & \textbf{Country} & \multicolumn{2}{l}{\textbf{Les Pays}} \\*
10.1 & Albania & Albanie & \ipa{albani} \\
10.2 & Austria & Autriche & \ipa{otʁiʃ} \\
10.3 & Belgium & Belgique & \ipa{bɛlʒik} \\
10.4 & Belorussia & Biélorussie & \ipa{bjelɔʁysi} \\
10.5 & Bosnia \& Herzegovina & Bosnie-Herzégovine & \ipa{bɔsni ɛʁzeɡɔvin} \\
10.6 & Bulgaria & Bulgarie & \ipa{bylɡaʁi} \\
10.7 & Croatia & Croatie & \ipa{kʁɔasi} \\
10.8 & Czech Republic & République Tchèque & \ipa{ʁepyblik tʃɛk} \\
10.9 & Denmark & Danemark & \ipa{danmaʁk} \\
10.11 & Eire & Irlande & \ipa{iʁlɑ̃d} \\
10.12 & England & Angleterre & \ipa{ɑ̃ɡlətɛʁ} \\
10.13 & Estonia & Estonie & \ipa{ɛstɔni} \\
10.14 & Finland & Finlande & \ipa{fɛ̃lɑ̃d} \\
10.15 & France & France & \ipa{fʁɑ̃s} \\
10.16 & Germany & Allemagne & \ipa{almaɲ} \\
10.17 & Gibraltar & Gibraltar & \ipa{ʒibʁaltaʁ} \\
10.18 & Greece & Grèce & \ipa{ɡʁɛs} \\
10.19 & Holland & Pays Bas / Hollande & \ipa{pe.i ba / ɔlɑ̃d} \\
10.21 & Hungary & Hongrie & \ipa{ɔ̃ɡʁi} \\
10.22 & Iceland & Islande & \ipa{islɑ̃d} \\
10.23 & Northern Ireland & Irlande du Nord & \ipa{iʁlɑ̃d dy nɔʁ} \\
10.24 & Isle of Man & Île de Man & \ipa{il də man} \\
10.25 & Italy & Italie & \ipa{itali} \\
10.26 & Latvia & Lettonie & \ipa{lɛtɔni} \\
10.27 & Liechtenstein & Liechtenstein & \ipa{liʃtɛnʃtajn} \\
10.28 & Lithuania & Lituanie & \ipa{litɥani} \\
10.29 & Luxembourg & Luxembourg & \ipa{lyksɑ̃buʁ} \\
10.31 & Macedonia & Macédoine & \ipa{masedwan} \\
10.32 & Moldavia & Moldavie & \ipa{mɔldavi} \\
10.33 & Netherlands & Pays Bas / Hollande & \ipa{pe.i ba / ɔlɑ̃d} \\
10.34 & Norway & Norvège & \ipa{nɔʁvɛʒ} \\
10.35 & Poland & Pologne & \ipa{pɔlɔɲ} \\
10.36 & Portugal & Portugal & \ipa{pɔʁtyɡal} \\
10.37 & Romania & Roumanie & \ipa{ʁumani} \\
10.38 & Russian Federation & Fédération Russe & \ipa{fedeʁasjɔ̃ ʁys} \\
10.39 & Scotland & Écosse & \ipa{ekɔs} \\
10.41 & Slovakia & Slovaquie & \ipa{slɔvaki} \\
10.42 & Slovenia & Slovénie & \ipa{slɔveni} \\
10.43 & Spain & Espagne & \ipa{ɛspaɲ} \\
10.44 & Svalbard & Svalbard (Spitzberg) & \ipa{svalbaʁ (spitsbɛʁɡ)} \\
10.45 & Sweden & Suède & \ipa{sɥɛd} \\
10.46 & Switzerland & Suisse & \ipa{sɥis} \\
10.47 & Turkey & Turquie & \ipa{tyʁki} \\
10.48 & Ukraine & Ukraine & \ipa{ykʁɛn} \\
10.49 & Wales & Pays de Galles & \ipa{pe.i də ɡal} \\
\end{longtable}
\normalsize
\section*{Hinweise zur Übertragung und Aussprache}
\small
\textbf{Vorlage und Umfang.} Grundlage ist „French v 1.0.doc“, mit den Credits
Jean-Michael -- F5OQO und Pierre-Louis -- F5BTT. Alle 142 Phrasen und Ländernamen sind enthalten.
Die ursprüngliche Gliederung einschließlich der Nummernsprünge und des fehlenden Abschnitts 7 bleibt erhalten.

\medskip
\textbf{Dokumentierte Korrekturen.} In 2.75 wurde das fehlende „de“ vor „l’Angleterre“ ergänzt,
in 2.78 das doppelte „au“ entfernt und in 8.5 „Qu'elle“ zu „Quelle“ berichtigt.
„Bosnie Herzégovie“ wurde zu „Bosnie-Herzégovine“ (10.5), „Ile de Man“ zu „Île de Man“ (10.24)
und „Ecosse“ zu „Écosse“ (10.39). „plait“ wurde als „plaît“ geschrieben, „Mhz“ als „MHz“
und „en Français“ als „en français“. Mehrfache Leerzeichen wurden vereinheitlicht;
das kaufmännische Und zwischen den Varianten in 9.18 wurde durch einen Zeilenumbruch ersetzt.

\medskip
\textbf{Abweichungen in der Vorlage.} Bei 2.2 steht im Englischen 1960, im Französischen 1930;
bei 4.6 stehen 3.558 MHz beziehungsweise 3,556 MHz. Beide Angaben bleiben jeweils unverändert.
Die Lautschrift folgt ausdrücklich der französischen Spalte: „mille neuf cent trente“ und
„trois virgule cinq cinq six mégahertz“. Peter / Pierre (2.1), „near mountains“ / „à la montagne“ (2.82)
und „great fun“ / „très intéressant“ (9.5) bleiben ebenfalls wie in der Vorlage.
Historische Ländernamen, Webadressen und die Formulierung „Ma carte QSL est sûre via bureau“
wurden nicht inhaltlich modernisiert. „Hollande“ wird bei 10.19 und 10.33 entsprechend der Vorlage als Alternative aufgeführt.

\medskip
\textbf{Transkriptionskonvention.} Die IPA-Spalte ist eine breite, manuell erstellte Aussprachehilfe
für Standardfranzösisch, keine enge Transkription einer Aufnahme. Nasalvokale sind mit einer Tilde markiert.
Die Unterscheidung \ipa{œ̃} / \ipa{ɛ̃} bleibt erhalten, obwohl viele Sprecher beide zusammenfallen lassen.
Unbetontes \ipa{ə} ist teilweise ausgeschrieben und kann im Gespräch entfallen. Französischer Wortgruppenakzent
wird nicht gesondert markiert. Der Punkt zeigt an ausgewählten Stellen eine Silbengrenze;
\ipa{‿} verdeutlicht Liaison oder eine Wortverbindung. Angegeben ist jeweils eine mögliche flüssige Aussprache;
optionale Liaisons und regionale Varianten sind nicht vollständig aufgeführt.

\medskip
\textbf{Namen, Kürzel und Zahlen.} Sittingbourne ist mit englischer, Hellschreiber mit deutscher Aussprache
wiedergegeben; sonst sind französische beziehungsweise französisch angepasste Formen verwendet.
Bei Markennamen, Ortsnamen und Fachwörtern sind Varianten möglich. Zahlen werden ausgeschrieben;
RST 599 wird als „R S T cinq neuf neuf“ gelesen. Rufzeichen und die meisten Kürzel werden mit französischen
Buchstabennamen gesprochen, nicht im internationalen Buchstabieralphabet. FISTS wird als Wort gelesen;
das e in eQSL und EHam als englisches \ipa{i}. Für CW steht „C double vé“, für RTTY „R T T i grec“.
8 °C wird als „huit degrés Celsius“ gesprochen. Klammern geben die optionalen Wörter der Vorlage wieder;
in 9.11 sind „souhaits“ und „vœux“ als zwei vollständige Varianten transkribiert.
Platzhalter \texttt{<\#\#\#\#\#>} sind nicht mitzusprechen und erscheinen in der Lautschrift als Auslassung.

\medskip
\textbf{LaTeX.} UTF-8-Datei; zweimal mit XeLaTeX kompilieren. Benötigt werden DejaVu Sans sowie die
üblichen LaTeX-Pakete der Präambel. Die IPA-Zeichen stehen direkt als Unicode im Quelltext.
Es sind keine externen Bilder oder Datendateien erforderlich.
\end{document}
